\section{The distance surface and its Brauer--Manin sets}
\label{sec:surface}

The four-distance equations also define a fixed surface, independent of
the square-free parameter in Section~\ref{sec:elliptic-sources}.
The surface approach is developed in \cite{berry1990}, and its
product-quotient description belongs to the general construction in
\cite{corvajaturchetzannier2025}.  We give the equations and their geometry
explicitly, and then determine the projection of the proper Brauer--Manin
set onto the positive real region.

\subsection{A canonical complete intersection}

Let $S\subset\PP^4_{[A:B:C:D:T]}$ be defined by
\begin{equation}\label{eq:surface-canonical}
 \begin{split}
 Q&=A^2+D^2-B^2-C^2=0,\\
 F&=(T^2+A^2-B^2)^2+(T^2+A^2-C^2)^2-4T^2A^2=0.
 \end{split}
\end{equation}
On $T=1$, the identification with the distance equations is
\begin{equation}\label{eq:surface-return}
 (r_0,r_1,r_2,r_3)=(A,B,C,D),\qquad
 x=\frac{1+A^2-B^2}{2},\qquad
 y=\frac{1+A^2-C^2}{2}.
\end{equation}
Indeed, $F=0$ gives $x^2+y^2=A^2$, and subtraction gives the other
three identities.  Let $X\to S$ be the minimal resolution.  Write $U$
for the open subset $T=1$, $ABCD\ne0$, viewed as an open subset of $X$,
and put
\[
 \mathcal C=\{(x,y,r_0,r_1,r_2,r_3)\in U(\R):
       0<x,y<1,\ r_j>0\text{ for all }j\}.
\]

\begin{proposition}\label{prop:surface-invariants}
The surface $S$ is geometrically integral and has precisely forty
ordinary double points.  Its resolution $X$ is a simply connected
minimal surface of general type, with
\[
 K_X^2=8,\qquad p_g(X)=5,\qquad q(X)=0,\qquad
 c_2(X)=64,\qquad b_2(X)=62.
\]
The geometric Picard group of $X$ is a finitely generated free abelian
group.
\end{proposition}

\begin{proof}
Work first over an algebraically closed field of characteristic zero.
Put $z=x+iy$, $w=x-iy$, and
\[
 (\alpha_0,\alpha_1,\alpha_2,\alpha_3)=(0,1,i,1+i).
\]
The affine distance equations become
\begin{equation}\label{eq:surface-branch-grid}
 r_j^2=(z-\alpha_j)(w-\overline\alpha_j),\qquad 0\le j\le3.
\end{equation}
Their four square classes are independent in
$\overline\Q(z,w)^*/\overline\Q(z,w)^{*2}$: the divisor
$z=\alpha_j$ occurs with odd multiplicity only in the $j$th radicand.
Thus the affine surface is geometrically integral.  The complete
intersection \eqref{eq:surface-canonical} is Cohen--Macaulay and purely
two-dimensional. Its intersection with $T=0$ has dimension one, so every
irreducible component meets the affine chart. There is consequently only
one component, and it is generically reduced. The Cohen--Macaulay property
excludes embedded components, hence $S$ is reduced and geometrically integral.

The branch grid in \eqref{eq:surface-branch-grid} determines all affine
singularities.  Away from the grid the covering is \'etale.  Along a
single branch line, its square root is a local parameter.  At an
intersection $z=\alpha_j$, $w=\overline\alpha_k$ with $j\ne k$, the two
corresponding square roots are independent local parameters.  At an
intersection with $j=k$, the local equation, after extracting square
roots of units, is $r_j^2=uv$.  This is an ordinary double point.
There are eight choices for the other three roots at each of the four
vertices, giving thirty-two affine nodes.  In distance coordinates the
three nonzero roots at a vertex are $\pm1,\pm1,\pm\sqrt2$.

At infinity, write $H_\pm=A^2-B^2\pm i(A^2-C^2)$.  The boundary is
the union of the two intersections $Q=H_\pm=0$ in $\PP^3$.  Each is
smooth: the ratios of the four diagonal coefficients of $H_\pm$ to
those of $Q$ are
\[
 1\pm i,\quad 1,\quad \pm i,\quad 0,
\]
which are pairwise distinct.  Proportional gradients would therefore
leave at most one coordinate nonzero, contrary to $Q=0$.  Their
intersections are exactly
\begin{equation}\label{eq:surface-infinity-nodes}
 [1:\epsilon_1:\epsilon_2:\epsilon_3:0],
 \qquad \epsilon_1,\epsilon_2,\epsilon_3\in\{1,-1\}.
\end{equation}
At any such point set $A=1$, $B=\epsilon_1+b$,
$C=\epsilon_2+c$, $T=t$, and eliminate $D$ using $Q$.
The quadratic term of $F$ is $4(b^2+c^2-t^2)$, which is nondegenerate.
These are eight more nodes.  The preceding smoothness calculation
also exhausts the points at infinity.  In particular $S$ is normal.

Adjunction for a complete intersection gives $\omega_S=\OO_S(1)$.
Ordinary double points are rational and their resolutions are crepant,
so
\[
 K_X=\pi^*\OO_S(1),\qquad K_X^2=\deg S=8.
\]
The Koszul resolution
\[
 0\longrightarrow\OO_{\PP^4}(-6)
 \longrightarrow\OO_{\PP^4}(-2)\oplus\OO_{\PP^4}(-4)
 \longrightarrow\OO_{\PP^4}
 \longrightarrow\OO_S\longrightarrow0
\]
gives $H^1(S,\OO_S)=0$ and $h^0(S,\OO_S(1))=5$.
Rationality of the singularities then gives $q=0$ and $p_g=5$ on $X$.
The canonical divisor is nef and big; adjunction excludes a
$(-1)$-curve.  Noether's formula gives $c_2=12(1-q+p_g)-K_X^2=64$,
and $b_1=2q=0$ gives $b_2=62$.

For simple connectedness, deform the two defining equations slightly
to a smooth complete intersection of type $(2,4)$, which is simply
connected by the Lefschetz hyperplane theorem.  Remove disjoint small
neighbourhoods of the forty nodes.  The smoothing and the resolution
have the same differentiable exterior.  Each boundary is
$\mathbf{RP}^3$.  Both the corresponding Milnor fibre and the
resolution neighbourhood are simply connected: the former retracts
onto a two-sphere, and the latter onto the exceptional projective
line.  Van Kampen's theorem therefore presents both fundamental
groups as the fundamental group of the common exterior modulo the
normal closures of the boundary-group images.  They are the same
quotient.  Thus $X(\mathbf C)$ is simply connected.

Consequently $H_1(X(\mathbf C),\Z)=0$, and the universal coefficient
theorem makes $H^2(X(\mathbf C),\Z)$ torsion-free.  The exponential
sequence and $q=0$ inject $\Pic(X_{\mathbf C})$ into this group.
Finally $\Pic^0(X)=0$ and the N\'eron--Severi theorem give finite
generation of the geometric Picard group.
\end{proof}

\subsection{A fixed product of genus-five curves}

Let $K=\Q(i)$.  Define $C_0/K$ by its function field
\begin{equation}\label{eq:surface-fixed-curve}
 K(C_0)=K(z)(u_0,u_1,u_2,u_3),\qquad u_j^2=z-\alpha_j,
\end{equation}
and let $\overline C_0$ denote its coefficient conjugate.  Both curves
are understood to be smooth and projective.

\begin{proposition}\label{prop:surface-product}
The curve $C_0$ has genus five.  The diagonal sign action of
$G=(\Z/2\Z)^4$ gives a birational identification
\[
 X_K\sim (C_0\times\overline C_0)/G.
\]
Its invariant functions include the four labelled products
$r_j=u_jv_j$, where $v_j^2=w-\overline\alpha_j$.
The Jacobian of $C_0$ is geometrically isogenous to a product of five
elliptic curves: four have $j=128$ and one has $j=1728$.
\end{proposition}

\begin{proof}
The four radicands in \eqref{eq:surface-fixed-curve} are independent
square classes by their valuations at their distinct zeros.  The
covering $C_0\to\PP^1$ has degree sixteen.  Its inertia elements are
the four individual sign changes at the finite branch points and their
sum at infinity.  Each inertia group has order two, so
\[
 2g(C_0)-2=16(-2)+5\cdot8=8.
\]
On the product, the diagonal action fixes $z,w,u_jv_j$.  The field
they generate has degree sixteen over $K(z,w)$ by the square-class
argument for \eqref{eq:surface-branch-grid}.  This is the degree of
the invariant field, since the product field has degree $256$ over
$K(z,w)$.  The invariant field is therefore precisely the distance
function field.

For a nonempty subset $J\subset\{0,1,2,3\}$, the corresponding
character quotient is
\[
 y_J^2=\prod_{j\in J}(z-\alpha_j).
\]
Only the four subsets of size three and the subset of size four
have positive genus; all five quotients have genus one.  The
pullbacks of their differentials give the five distinct character
lines in $H^0(C_0,\Omega^1)$, hence their product induces an isogeny
on Jacobians.  The branch cross-ratios can be chosen as $i$ for the
three-element quotients and $2$ for the four-element quotient.
Substitution in
\[
 j(\lambda)=256\frac{(1-\lambda+\lambda^2)^3}
                         {\lambda^2(1-\lambda)^2}
\]
gives the asserted values.
\end{proof}

\subsection{The positive real component}

We write $\operatorname{Br}$ for the cohomological Brauer group and
$\operatorname{Br}_0$ for the image of the Brauer group of the ground
field.  The rational point used below lies on the boundary
of $U$.

\begin{lemma}\label{lem:surface-real-boundary}
There is a point $P_*\in X(\Q)$ in the same real connected
component as every point of $\mathcal C$.
\end{lemma}

\begin{proof}
At the node $[1:1:1:1:0]$, use the local coordinates $b,c,t$ from
the proof of Proposition~\ref{prop:surface-invariants}.  The
exceptional conic is $b^2+c^2=t^2$.  Take the rational point
$P_*=[-1:0:1]$ on this conic.

Start with a point of $\mathcal C$, fix its $y\in(0,1)$, and let
$x$ tend to $+\infty$, choosing all four positive distance roots.
This is a path in the smooth affine locus because none of the
distances vanishes.  In the chart $A=1$ its coordinates are
$B=r_1/r_0$, $C=r_2/r_0$, $T=1/r_0$.  Rationalization gives
\[
 \frac b t=r_1-r_0=\frac{1-2x}{r_1+r_0}\longrightarrow-1,
 \qquad
 \frac c t=r_2-r_0=\frac{1-2y}{r_2+r_0}\longrightarrow0.
\]
The lifted path extends to $P_*$ on the resolution.
\end{proof}

\begin{proposition}\label{prop:surface-brauer-finite}
The group $\operatorname{Br}(X)/\operatorname{Br}_0(X)$ is finite.
\end{proposition}

\begin{proof}
Resolve the rational map $C_0\times\overline C_0\dashrightarrow X_K$
from Proposition~\ref{prop:surface-product}.  This gives a smooth
projective surface $Y$, birational to the product, and a generically
finite morphism $f:Y\to X_K$ of degree sixteen.
The product theorem of Skorobogatov--Zarhin
\cite[Theorem~A]{skorobogatovzarhin2014}, with its corrected proof
\cite{skorobogatovzarhin2023}, states that, over a field finitely
generated over $\Q$, the map
\[
 \operatorname{Br}(\overline V)^{\Gamma}\oplus
 \operatorname{Br}(\overline W)^{\Gamma}
 \longrightarrow
 \operatorname{Br}(\overline V\times\overline W)^{\Gamma}
\]
has finite cokernel for smooth projective geometrically integral
varieties $V,W$.  The geometric Brauer group of a smooth projective
curve is zero.  Applied to our two curves, this theorem and
birational invariance therefore show that
$\operatorname{Br}(\overline Y)^{\Gamma_K}$ is finite.

Restriction--corestriction on function fields shows that the kernel of
$f^*: \operatorname{Br}(\overline X)\to\operatorname{Br}(\overline Y)$
is killed by sixteen.  This kernel is finite, because the Kummer
sequence places $\operatorname{Br}(\overline X)[16]$ in a quotient of
the finite group $H^2_{\mathrm{\acute et}}(\overline X,\mu_{16})$.
Taking invariant subgroups thus proves that
$\operatorname{Br}(\overline X)^{\Gamma_K}$, and hence
$\operatorname{Br}(\overline X)^{\Gamma_\Q}$, is finite.

By Proposition~\ref{prop:surface-invariants}, $\Pic(\overline X)$
is a finitely generated free abelian group.  Its Galois action factors
through a finite group.  Inflation--restriction and the vanishing of
continuous homomorphisms from a profinite group to a discrete free
abelian group identify its first Galois cohomology with that of a
finite group acting on a lattice; this group is finite.  The
Hochschild--Serre sequence now embeds
$\operatorname{Br}_1(X)/\operatorname{Br}_0(X)$ in
$H^1(\Q,\Pic(\overline X))$, while
$\operatorname{Br}(X)/\operatorname{Br}_1(X)$ embeds in
$\operatorname{Br}(\overline X)^{\Gamma_\Q}$.  Both are finite.
\end{proof}

\subsection{Adelic points with a prescribed strict real component}

For a subgroup $B\subset\operatorname{Br}(U)$, let
$U(\mathbf A_\Q)^B$ denote the points in the usual restricted-product
adelic space orthogonal to $B$ under the Brauer--Manin pairing.  We
identify $\operatorname{Br}(X)$ with its restriction to $U$.

\begin{theorem}\label{thm:surface-brauer-projection}
Every $w\in\mathcal C$ is the real component of an adelic point in
$U(\mathbf A_\Q)^{\operatorname{Br}(X)}$.  Moreover, for every finite
subgroup $B\subset\operatorname{Br}(U)$ and every $w\in\mathcal C$,
there is a point $(P_v)\in U(\mathbf A_\Q)^B$ with $P_\infty=w$.
\end{theorem}

\begin{proof}
First let $B_0$ be a finite subgroup of $\operatorname{Br}(X)$.
Choose a finite set $\Sigma$ of primes such that $X$, $U$ and
$P_*$ extend to a smooth proper model $\mathscr X$, a smooth
geometrically integral open model $\mathscr U$, and a section over
$\Z[1/\Sigma]$, and all the classes in $B_0$ extend to
$\operatorname{Br}(\mathscr X)$.  The section is denoted by the same
symbol $P_*$ here.  Lang--Weil and Hensel's lemma, after enlarging
$\Sigma$, give a point of $\mathscr U(\Z_p)$ for every
$p\notin\Sigma$.  Every class in $B_0$ evaluates to zero at this
point and at the section, since $\operatorname{Br}(\Z_p)=0$.

At each $p\in\Sigma$, the evaluation of a Brauer class on $X(\Q_p)$
is locally constant.  Choose a sufficiently small neighbourhood of
the rational point $P_*$ on which all the evaluations in $B_0$
agree with their values at $P_*$.  This neighbourhood contains a
point of $U(\Q_p)$: in smooth analytic coordinates a proper algebraic
closed subset cannot contain an open ball.  At the real place use
the prescribed $w$.  Evaluation is constant on real connected
components, so Lemma~\ref{lem:surface-real-boundary} gives agreement
with $P_*$ there as well.

These choices form an adelic point of $U$, integral in the single
model $\mathscr U$ outside $\Sigma$.  For each $\beta\in B_0$ its
evaluations coincide place by place with those of
$\beta(P_*)\in\operatorname{Br}(\Q)$.  Global reciprocity
therefore gives orthogonality to $B_0$.

Proposition~\ref{prop:surface-brauer-finite} supplies finitely many
representatives for
$\operatorname{Br}(X)/\operatorname{Br}_0(X)$.  Because the Brauer
group is torsion, they generate a finite subgroup $B_0$.  Applying
the construction to this group gives orthogonality to all of
$\operatorname{Br}(X)$, as constants have zero total invariant.

For the last assertion, use Harari's formal lemma in the form
recalled in \cite{colliothelene_skoro2021}: if $V$ is smooth and
geometrically integral over a number field, $U\subset V$ is a
nonempty open subset, and $B\subset\operatorname{Br}(U)$ is finite,
then an adelic point of $U$ orthogonal to
$B\cap\operatorname{Br}(V)$ can be replaced by one orthogonal to
$B$, preserving its components at any prescribed finite set of
places.  Apply the lemma with $V=X$, starting with the point just
constructed, and preserve the real component.
\end{proof}

Thus the positive real region lies in the projection of the adelic
set for the entire proper Brauer group, and also in the projection
for each finite subgroup of the affine Brauer group.  In the second
assertion the finite subgroup is fixed before the adelic point is
chosen; this is the quantifier order used in the theorem.
